\section{Problem Statement, Main Concepts and Implementation Strategies}

As mentioned at the beginning of this work, programming languages must constantly evolve their capabilities and adapt to the technological and economic demands of their era. However, they must do so while preserving their core identity---their language architecture, fundamental structures, memory management model, and built-in primitives. As we established early on, there is no perfect programming language. A language's effectiveness cannot simply be attributed to the wisdom of its creators; rather, a language emerges to solve specific problems within a particular temporal, resource, and technological context. Language designers cannot fully anticipate future technological advances or societal needs.

Given these considerations, and recognizing that artificial intelligence, data analysis, and high-performance computing represent the most significant technological drivers of our time, we undertook the challenge of creating an environment for developing heterogeneous applications using TypeScript. Our decision was not born from spontaneous impulse or a desire for recognition in the programming community. Rather, guided by the Red Queen hypothesis, we recognized that such transformation is inevitable if the TypeScript/JavaScript ecosystem wishes to remain competitive and relevant.

TypeScript/JavaScript has already demonstrated its ability to meet diverse challenges. Technologies like NodeJS, Bun, and Deno exemplify this adaptability in server-side computing. Furthermore, frameworks such as ReactJS, React Native, Babel, and NextJS have proven TypeScript/JavaScript to be an exceptionally flexible language for application development. Libraries like numeric-js, math-js, and d3 show promise for computational tasks, though the ecosystem currently lacks the maturity of established \acr{HPC} languages.

This observation leads to a crucial insight: while Python currently dominates high-performance computing through mature libraries like NumPy, SciPy, TensorFlow, and PyTorch (which leverage optimized C/Fortran backends), TypeScript possesses several \textit{latent advantages} that could make it competitive with proper tooling. TypeScript offers superior baseline performance compared to Python, native type support that facilitates optimization, and versatility across web, mobile, and server applications. More importantly, TypeScript's well-defined grammar and the availability of robust parsing infrastructure (through the TypeScript compiler \acr{API}) provide an excellent foundation for building advanced compilation toolchains.

The key limitation is not the language itself, but rather the absence of infrastructure to compile TypeScript to low-level, hardware-optimized code. This is precisely the gap we aim to bridge: by leveraging \acr{LLVM} \acr{IR} as a compilation target, we can transform TypeScript from a primarily application-development language into a viable platform for heterogeneous computing, including \acr{GPU} acceleration and multi-architecture deployment. Our approach does not claim that TypeScript is inherently superior for \acr{HPC}; rather, we demonstrate that with appropriate compiler infrastructure, TypeScript \textit{can become} a practical choice for high-performance computing workloads.


Existing libraries and technologies that support heterogeneous programming follow the foreign function interface (\acr{FFI}) pattern: high-performance code written in languages suited for computational workloads---such as C, Fortran, or \acr{CUDA}---is invoked from TypeScript, mirroring established practices in Python, R, and MATLAB. While this approach has merit, its fundamental limitation is the requirement for developers to possess expertise across multiple programming languages. The TypeScript ecosystem does offer such technologies, including \texttt{cuda-ts} and \texttt{llvm-bindings}, which enable developers to execute \acr{LLVM} \acr{IR} and \acr{CUDA} code. However, these solutions still require developers to write their computational kernels separately in \acr{CUDA} or \acr{LLVM} \acr{IR}, maintaining the multi-language burden.

The TypeScript ecosystem has begun to show early signs of native heterogeneous programming capabilities. Libraries such as \texttt{GPU.js} and \texttt{taichi.js} represent key developments in this direction, transpiling JavaScript/TypeScript to WebGL or WebGPU instructions to achieve \acr{GPU} acceleration within browser environments. Similarly, TensorFlow.js leverages WebGPU-based acceleration to significantly improve performance for neural network operations. However, solutions based on WebGL or WebGPU suffer from the inherent constraints of browser graphics \acr{API}s. They cannot exploit \acr{NVIDIA}-specific features such as shared memory hierarchies or warp-level primitives, nor can they match the performance of native \acr{CUDA} kernels for general-purpose \acr{GPU} computing. Another technology, AssemblyScript---a TypeScript-like language---compiles to WebAssembly for \acr{CPU} execution but lacks \acr{GPU} code generation capabilities entirely. 


Thus, our objective in this chapter is to \textit{\textbf{build a compiler from TypeScript to \acr{LLVM} \acr{IR} and to \acr{NVPTX} and \acr{AMDGPU} targets}}. We have already outlined the fundamental steps of compiler construction in Section~\ref{subsec:compiler-phases}. Here, we will examine what has been achieved in this highly specialized domain, providing guidance and insights for accomplishing our stated goal.

We will begin by exploring the TypeScript compiler \acr{API} documentation, as it enables us to complete the initial phases of the compilation pipeline previously described. Subsequently, we will employ our \texttt{@euriklis/llvm-ir} library to generate the necessary \acr{LLVM} \acr{IR} for the target platform from the TypeScript \acr{AST} (corresponding to step 4 in Section~\ref{subsec:compiler-phases}). During this phase, we will also leverage the TypeScript transpilation \acr{API} to transform the source code and identify the classes containing methods that require compilation to \acr{GPU} kernels.

This methodology draws inspiration from the Python Numba project. However, we cannot replicate Numba's implementation exactly in TypeScript. This constraint presents both advantages and disadvantages, though fortunately the advantages are more numerous and significant.

The primary limitation is that we cannot reproduce Numba's exact structure and methods in identical form within TypeScript. This inability to port the technology's surface-level \acr{API} stems from fundamental differences between the two languages. Specifically, while both languages support decorators, TypeScript decorators cannot be applied to standalone functions as they can in Python---they only work with class members. This restriction posed a significant design challenge. To address this limitation, we introduce a \texttt{HeterogeneousProgram} class that serves as a container for methods requiring \acr{GPU} compilation.

Beyond the decorator limitation, several other disadvantages emerge when adapting Numba's approach to TypeScript. TypeScript's numeric computing ecosystem lacks the maturity of Python's scientific stack---there is no direct equivalent to NumPy's optimized array operations or SciPy's comprehensive numerical algorithms that Numba seamlessly integrates with. Additionally, TypeScript's type system, while providing excellent compile-time guarantees, suffers from \textit{type erasure}: types are stripped away during transpilation to JavaScript, requiring our compiler to perform separate type analysis during the compilation pipeline. The \acr{HPC} community around TypeScript is also considerably smaller, meaning fewer reference implementations, less accumulated expertise in \acr{GPU} computing patterns, and a more limited ecosystem of interoperable libraries. Finally, TypeScript lacks Python's direct C/Fortran interoperability mechanisms (ctypes, Cython, CFFI), which Numba exploits to interface with existing high-performance libraries.

However, TypeScript offers several compelling advantages that motivated this project. Most significantly, TypeScript's \textit{static type system} provides type information at compile time through explicit annotations, eliminating the need for complex runtime type inference that Numba must perform on Python code. This compile-time type information enables more aggressive optimizations and simplifies the compiler implementation. TypeScript also benefits from exceptional tooling infrastructure: the TypeScript compiler exposes a comprehensive \acr{API} for \acr{AST} manipulation, transformation, and type checking, providing a robust foundation for building compilation toolchains. The language's position within the JavaScript ecosystem opens unique opportunities---our compiler could potentially enable \acr{GPU}-accelerated computing in browser environments via WebGPU, bridging high-performance computing with web applications in ways that Python cannot. Unlike Python's Global Interpreter Lock (\acr{GIL}), TypeScript has a simpler runtime model that avoids concurrency bottlenecks, and modern JavaScript engines (V8, JavaScriptCore) already deliver superior baseline performance compared to CPython. Furthermore, TypeScript's cross-platform nature allows the same codebase to target server, desktop, and client environments, offering deployment flexibility that traditional \acr{HPC} languages lack. These advantages---particularly the static types, excellent compiler \acr{API}, and potential for browser-based \acr{GPU} computing---make TypeScript a compelling platform for heterogeneous computing despite the challenges of diverging from Numba's established patterns. 
